% TODO proof-reading
\subsection{ループ}

\begin{lstlisting}[style=customjava]
public class Loop
{
	public static void main(String[] args)
	{ 
		for (int i = 1; i <= 10; i++)
		{
			System.out.println(i); 
		}               
	}
}
\end{lstlisting}

\begin{lstlisting}
  public static void main(java.lang.String[]);
    flags: ACC_PUBLIC, ACC_STATIC
    Code:
      stack=2, locals=2, args_size=1
         0: iconst_1      
         1: istore_1      
         2: iload_1       
         3: bipush        10
         5: if_icmpgt     21
         8: getstatic     #2         // Field java/lang/System.out:Ljava/io/PrintStream;
        11: iload_1       
        12: invokevirtual #3         // Method java/io/PrintStream.println:(I)V
        15: iinc          1, 1
        18: goto          2
        21: return        
\end{lstlisting}

\TT{iconst\_1}は1を\ac{TOS}にロードし、\TT{istore\_1}はそれを\ac{LVA}のスロット1に格納します。

なぜ0番目のスロットではないのでしょうか？
\main関数には1つの引数（\TT{String}の配列）があり、
それへのポインタ（または\IT{reference}）が現在0番目のスロットにあるからです。

したがって、\IT{i}ローカル変数は常に1番目のスロットにあります。

オフセット3と5の命令は\IT{i}と10を比較します。

\IT{i}が大きい場合、実行フローは関数が終了するオフセット21に移ります。

そうでなければ、\TT{println}が呼び出されます。

\IT{i}はその後、\TT{println}のためにオフセット11で再ロードされます。

ちなみに、\IT{integer}用の\TT{println}メソッドを呼び出しており、
これをコメントで見ることができます：「(I)V」
（\IT{I}は\IT{integer}を意味し、\IT{V}は戻り型が\IT{void}であることを意味します）。

\TT{println}が終了すると、\IT{i}はオフセット15でインクリメントされます。

命令の最初のオペランドはスロット番号（1）で、
2番目は変数に加算する数（1）です。

\TT{goto}は単なるGOTOで、ループ本体の開始であるオフセット2にジャンプします。

より複雑な例に進みましょう：

\begin{lstlisting}[style=customjava]
public class Fibonacci
{
	public static void main(String[] args)
	{ 
		int limit = 20, f = 0, g = 1;

		for (int i = 1; i <= limit; i++)
		{
			f = f + g;
			g = f - g;
			System.out.println(f); 
		}
	}
}
\end{lstlisting}

\begin{lstlisting}
  public static void main(java.lang.String[]);
    flags: ACC_PUBLIC, ACC_STATIC
    Code:
      stack=2, locals=5, args_size=1
         0: bipush        20
         2: istore_1      
         3: iconst_0      
         4: istore_2      
         5: iconst_1      
         6: istore_3      
         7: iconst_1      
         8: istore        4
        10: iload         4
        12: iload_1       
        13: if_icmpgt     37
        16: iload_2       
        17: iload_3       
        18: iadd          
        19: istore_2      
        20: iload_2       
        21: iload_3       
        22: isub          
        23: istore_3      
        24: getstatic     #2         // Field java/lang/System.out:Ljava/io/PrintStream;
        27: iload_2       
        28: invokevirtual #3         // Method java/io/PrintStream.println:(I)V
        31: iinc          4, 1
        34: goto          10
        37: return        
\end{lstlisting}
        
これは\ac{LVA}スロットのマップです：

\begin{itemize}
\item 0 --- \mainの唯一の引数
\item 1 --- \IT{limit}、常に20を含む
\item 2 --- \IT{f}
\item 3 --- \IT{g}
\item 4 --- \IT{i}
\end{itemize}

Javaコンパイラが、ソースコードで宣言された順序と同じ順序で\ac{LVA}スロットに変数を割り当てることがわかります。

スロット0、1、2、3にアクセスするための個別の\TT{istore}命令がありますが、
4以上のスロットにはありません。そのため、オフセット8にはスロット番号をオペランドとして取る
追加のオペランドを持つ\TT{istore}があります。

オフセット10の\TT{iload}でも同じです。

しかし、常に20を含む（本質的に定数である）\IT{limit}変数に別のスロットを割り当て、
その値を頻繁に再ロードするのは疑わしいのではないでしょうか？

\ac{JVM} \ac{JIT}コンパイラは通常、そのようなことを最適化するのに十分優秀です。

コードへの手動介入はおそらく価値がありません。